\SourceLeaf{252}der Brennweite, welche die meisten Schriftsteller als die Entfernung des Brennpunkts der Linse von ihrem Mittelpunkte erklären, indem sie von vorne her entweder stillschweigend voraussetzen, oder ausdrücklich bevorworten, dass die Dicke der Linse hiebei wie unendlich klein betrachtet werde, wodurch also für wirkliche Linsen die Brennweite eine Unbestimmtheit von der Ordnung der Dicke der Linsen behält. Wo es einmal genauer genommen wird, rechnet man jene Entfernung bald von der dem Brennpunkte nächsten Oberfläche der Linse, bald von dem sogenannten optischen Mittelpunkte derselben, bald von demjenigen Punkte, welcher zwischen der Vorderfläche und Hinterfläche mitten inne liegt, und von allen diesen Bestimmungen wieder verschieden ist derjenige Werth, welcher bei der Vergleichung der Grösse des Bildes eines unendlich entfernten Gegenstandes mit der scheinbaren Grösse des letztern zum Grunde gelegt werden muss, welche letztere Bestimmung in der That die einzige zweckmässige ist.

Ich habe daher für nicht überflüssig gehalten, diesen an sich ganz elementaren Untersuchungen einige Blätter zu widmen, vornehmlich um zu zeigen, dass bei den oben erwähnten eleganten Sätzen ohne Verlust für ihre Einfachheit die Dicke der Linsen mit berücksichtigt werden kann. Nur die Beschränkung auf solche Strahlen, die gegen die Axe unendlich wenig geneigt sind, soll hier beibehalten, oder die Abweichung wegen der Kugelgestalt bei Seite gesetzt werden.

\Article{1}
Die Bestimmung der Lage aller in dieser Untersuchung vorkommenden Punkte geschieht durch rechtwinklige Coordinaten $x,y,z$, wobei vorausgesetzt wird, dass die Mittelpunkte der verschiedenen Brechungsflächen in der Axe der $x$ liegen, und nur solche Lichtstrahlen betrachtet werden, die mit dieser Axe einen sehr kleinen Winkel machen: die Coordinaten $x$ werden, bei ganz willkürlichem Anfangspunkte, als wachsend angenommen in dem Sinne der Richtung der Lichtstrahlen.

Wir betrachten zuerst die Wirkung Einer Brechung auf den Weg eines Lichtstrahls. Es sei das Brechungsverhältniss beim Übergange aus dem ersten Mittel in das zweite wie $\frac1n$ zu $\frac1{n'}$, oder wie $n'$ zu $n$. Wir bezeichnen mit $M$ den Mittelpunkt der sphärischen Scheidungsfläche zwischen den beiden Mitteln, mit $N$ den Durchschnittspunkt dieser Fläche mit der ersten Coordinatenaxe; zugleich sollen mit denselben Buchstaben auch die diesen Punkten entsprechenden
\SourceLeaf{253}Werthe von $x$ bezeichnet werden, was in der Folge auch bei andern Punkten der ersten Coordinatenaxe eben so gehalten werden soll. Es sei ferner $r=M-N$, oder $r$ der Halbmesser der Scheidungsfläche, positiv oder negativ, je nachdem das erste Mittel an der convexen oder an der concaven Seite liegt; $P$ der Punkt, wo der Lichtstrahl die Scheidungsfläche trifft, und $\theta$ der (spitze) Winkel zwischen $MP$ und der Axe der $x$.

Die von einem Lichtstrahle vor der Brechung beschriebene gerade Linie wird durch zwei Gleichungen bestimmt, denen wir folgende Formel geben:
\[
\begin{aligned}
y&=\frac\beta n(x-N)+b\\
z&=\frac\gamma n(x-N)+c
\end{aligned}
\]
und eben so seien die Gleichungen für die von demselben Lichtstrahle nach der Brechung beschriebene gerade Linie
\[
\begin{aligned}
y&=\frac{\beta'}{n'}(x-N)+b'\\
z&=\frac{\gamma'}{n'}(x-N)+c'
\end{aligned}
\]
Es kommt also darauf an, die Abhängigkeit der vier Grössen $\beta',\gamma',b',c'$ von $\beta,\gamma,b,c$ zu entwickeln. Für den Punkt $P$ wird
\[
x=N+r(1-\cos\theta)
\]
also, weil für denselben sowohl die ersten als die zweiten Gleichungen gelten,
\[
\frac\beta n.r(1-\cos\theta)+b=\frac{\beta'}{n'}.r(1-\cos\theta)+b'
\]
und folglich, da $\beta,\beta',\theta$ als unendlich kleine Grössen erster Ordnung gelten, bis auf Grössen dritter Ordnung genau
\[
b'=b\quad .\quad .\quad .\quad .\quad .\quad .\quad .\quad .\quad .\quad .\quad .\qquad(1)
\]
und eben so
\[
c'=c\quad .\quad .\quad .\quad .\quad .\quad .\quad .\quad .\quad .\quad .\quad .\qquad(1)
\]

Eine durch $M$ senkrecht gegen die Axe der $x$ gelegte Ebene werde von dem ersten (nöthigenfalls verlängerten) Wege des Lichtstrahls in $Q$, von dem zweiten in $Q'$ geschnitten. Da $PQ'$ mit $PQ$ und $PM$ in Einer Ebene liegt, so sind $M,Q,Q'$ in Einer geraden Linie. Bezeichnet man mit $\lambda,\lambda'$ die Winkel, welche diese gerade Linie mit $PQ,PQ'$ macht, so werden offenbar
\SourceLeaf{254}$MQ.\sin\lambda,MQ'.\sin\lambda'$ den Producten aus dem positiv genommenen Halbmesser der Kugelfläche in die Sinus des Einfallswinkels und des gebrochenen Winkels gleich, also den Zahlen $n',n$ proportional sein, mithin
\[
MQ'=\frac{n.MQ.\sin\lambda}{n'\sin\lambda'}
\]
Da nun für den Punkt $Q$
\[
\begin{aligned}
y&=b+\frac{\beta r}n\\
z&=c+\frac{\gamma r}n
\end{aligned}
\]
für den Punkt $Q'$ hingegen
\[
\begin{aligned}
y&=b'+\frac{\beta'r}{n'}\\
z&=c'+\frac{\gamma'r}{n'}
\end{aligned}
\]
wird, und die beiden letztern Coordinaten sich zu den beiden erstern wie $MQ'$ zu $MQ$ verhalten, so hat man
\[
\begin{aligned}
b'+\frac{\beta'r}{n'}&=\frac{n\sin\lambda}{n'\sin\lambda'}.\left(b+\frac{\beta r}n\right)\\
c'+\frac{\gamma'r}{n'}&=\frac{n\sin\lambda}{n'\sin\lambda'}.\left(c+\frac{\gamma r}n\right)
\end{aligned}
\]
oder
\[
\begin{aligned}
\beta'&=\frac{nb+\beta r}{r}.\frac{\sin\lambda}{\sin\lambda'}-\frac{n'b'}r\\
\gamma'&=\frac{nc+\gamma r}{r}\raisebox{0.55em}{.}\frac{\sin\lambda}{\sin\lambda'}-\frac{n'c'}r
\end{aligned}
\]
Diese Ausdrücke sind strenge richtig; allein, da $\lambda,\lambda'$ vom rechten Winkel um Grössen erster Ordnung, also ihre Sinus von der Einheit um Grössen zweiter Ordnung verschieden sind, so wird, auf Grössen dritter Ordnung genau,
\[
\left.\begin{aligned}
\beta'&=\beta-\frac{n'-n}r.b=\beta+\frac{n'-n}{N-M}.b\\
\gamma'&=\gamma-\frac{n'-n}r.c=\gamma+\frac{n'-n}{N-M}.c
\end{aligned}\right\}\quad .\quad .\quad .\quad .\quad .\quad .\qquad(2)
\]
Diese Gleichungen (1), (2) enthalten die Auflösung unsrer Aufgabe.

Es verdient bemerkt zu werden, dass dieselben Formeln auch unmittelbar auf einen zurückgeworfenen Strahl angewandt werden können, wenn man nur $-n$ für $n'$ substituirt, und dass, mit Hülfe eines solchen Verfahrens, auch die sämmtlichen folgenden Untersuchungen sich sehr leicht auf den Fall erweitern lassen, wo anstatt der Refractionen eine oder mehrere Reflexionen eintreten.

\SourceLeaf{255}\Article{2}
Zur Auflösung der allgemeinen Aufgabe, den Weg des Lichtstrahls nach wiederholten $(\mu+1)$ Brechungen zu bestimmen, wollen wir folgende Bezeichnungen gebrauchen.

$N^0,N',N''\, .\, .\, .\, .\, .\, N^{(\mu)}$ die Punkte, wo die Axe der $x$ von den Brechungsflächen getroffen wird.

$M^0,M',M''\, .\, .\, .\, .\, M^{(\mu)}$ die in dieser Axe liegenden Mittelpunkte der Brechungsflächen

$n':n^0,\ n'':n',\ n''':n''\, .\, .\, .\, .\, .\, n^{(\mu+1)}:n^{(\mu)}$ die Brechungsverhältnisse beim Durchgange aus dem ersten Mittel (vor $N^0$) in das zweite (zwischen $N^0$ und $N'$), aus dem zweiten in das dritte u.\,s.\,f.\ In der Emanationstheorie sind also die Zahlen $n^0,n',n''$ u.\,s.\,w den Geschwindigkeiten der Fortpflanzung des Lichts in den einzelnen Mitteln direct, in der Undulationstheorie verkehrt proportional, und wenn das letzte Mittel dasselbe ist, wie das erste, wird $n^{(\mu+1)}=n^0$.

Die Gleichungen für den Weg des Lichtstrahls vor der ersten Brechung seien
\[
\begin{aligned}
y&=\frac{\beta^0}{n^0}(x-N^0)+b^0\\
z&=\frac{\gamma^0}{n^0}(x-N^0)+c^0
\end{aligned}
\]
die Gleichungen für den Weg nach der ersten Brechung folgende
\[
\begin{aligned}
y&=\frac{\beta'}{n'}(x-N^0)+b^0\\
z&=\frac{\gamma'}{n'}(x-N^0)+c^0
\end{aligned}
\]
oder, anstatt auf $N^0$, auf $N'$ bezogen
\[
\begin{aligned}
y&=\frac{\beta'}{n'}(x-N')+b'\\
z&=\frac{\gamma'}{n'}(x-N')+c'
\end{aligned}
\]
eben so die Gleichungen für den Weg nach der zweiten Brechung
\[
\begin{aligned}
y&=\frac{\beta''}{n''}(x-N')+b'\\
z&=\frac{\gamma''}{n''}(x-N')+c'
\end{aligned}
\]
oder
\[
\begin{aligned}
y&=\frac{\beta''}{n''}(x-N'')+b''\\
z&=\frac{\gamma''}{n''}(x-N'')+c''
\end{aligned}
\]
\SourceLeaf{256}u. s. f., also, wenn wir die letzten Glieder in den Reihen der $\beta,\gamma,n,N,b,c$, nemlich $\beta^{(\mu+1)},\gamma^{(\mu+1)},n^{(\mu+1)},N^{(\mu)},b^{(\mu)},c^{(\mu)}$, um sie als solche kenntlich zu machen, durch $\beta^*,\gamma^*,n^*,N^*,b^*,c^*$ bezeichnen, die Gleichungen für den letzten Weg des Lichtstrahls
\[
\begin{aligned}
y&=\frac{\beta^*}{n^*}(x-N^*)+b^*\\
z&=\frac{\gamma^*}{n^*}(x-N^*)+c^*
\end{aligned}
\]
Endlich setzen wir zur Abkürzung
\[
\left.\begin{aligned}
\frac{N'-N^0}{n'}&=t',&\frac{N''-N'}{n''}&=t'',&\frac{N'''-N''}{n'''}&=t'''\ \text{u. s. f.}\\
\frac{n'-n^0}{N^0-M^0}&=u^0,&\frac{n''-n'}{N'-M'}&=u',&\frac{n'''-n''}{N''-M''}&=u''\ \text{u. s. f.}
\end{aligned}\right\}\quad .\quad .\qquad (3)
\]
und der Analogie nach für die letzten Glieder in diesen Reihen
\[
t^{(\mu)}=t^*,\qquad u^{(\mu)}=u^*
\]
Es wird demnach, in Folge des vorhergehenden Artikels,
\[
\begin{aligned}
\beta'&=\beta^0+u^0b^0\\
b'&=b^0+t'\beta'\\
\beta''&=\beta'+u'b'\\
b''&=b'+t''\beta''\\
\beta'''&=\beta''+u''b''\\
b'''&=b''+t'''\beta'''
\end{aligned}
\]
u. s. f., woraus erhellt, dass $b^*,\beta^*$ linearisch durch $b^0$ und $\beta^0$ bestimmt werden, und dass, wenn man
\[
\left.\begin{aligned}
b^*&=gb^0+h\beta^0\\
\beta^*&=kb^0+l\beta^0
\end{aligned}\right\}\quad .\quad .\quad .\quad .\quad .\quad .\quad .\quad .\quad .\qquad (4)
\]
setzt, in der von \textsc{Euler} (\textit{Comment. Nov. Acad. Petropol.} T. IX) eingeführten Bezeichnung sein wird
\[
\left.\begin{aligned}
g&=(u^0,t',u',t'',u''\ .\ .\ .\ .\ .\ t^*)\\
h&=(t',u',t'',u''\ .\ .\ .\ .\ .\ t^*)\\
k&=(u^0,t',u',t'',u''\ .\ .\ .\ .\ .\ .\ u^*)\\
l&=(t',u',t'',u''\ .\ .\ .\ .\ .\ .\ u^*)
\end{aligned}\right\}\quad .\quad .\quad .\quad .\quad .\quad .\quad .\qquad (5)
\]
\SourceLeaf{257}Die Bedeutung dieser Bezeichnung besteht bekanntlich darin, dass, wenn aus einer gegebenen Reihe von Grössen $a,a',a'',a'''$ u. s. f. eine andere Reihe, $A,A',A'',A'''$ u. s. f. nach folgendem Algorithmus gebildet wird
\[
A=a,\quad A'=a'A+1,\quad A''=a''A'+A,\quad A'''=a'''A''+A'\quad\text{u. s. f.}
\]
man schreibt
\[
A=(a),\quad A'=(a,a'),\quad A''=(a,a',a''),\quad A'''=(a,a',a'',a''')\quad\text{u. s. f.}
\]
Übrigens ist von selbst klar, dass in den Gleichungen für die dritte Coordinate $z$ die Constanten für den letzten Weg aus denen für den ersten ganz eben so abgeleitet werden, wie in den Gleichungen für $y$, oder dass man haben wird
\[
\left.\begin{aligned}
c^*&=gc^0+h\gamma^0\\
\gamma^*&=kc^0+l\gamma^0
\end{aligned}\right\}\quad .\quad .\quad .\quad .\quad .\quad .\quad .\quad .\quad .\qquad (4)
\]
In den Gleichungen (3), (5), (4) ist die vollständige Auflösung unsrer Aufgabe enthalten.

\Article{3}
\textsc{Euler} hat a. a. O. die vornehmsten den erwähnten Algorithmus betreffenden Relationen entwickelt, von denen hier nur zwei in Erinnerung gebracht werden mögen.

Erstlich ist immer
\[
(a,a',a''\ldots a^{(\lambda)})(a',a''\ldots a^{(\lambda+1)})-(a,a',a''\ldots a^{(\lambda+1)})(a',a''\ldots a^{(\lambda)})=\pm1
\]
wo das obere oder das untere Zeichen gilt, je nachdem die Anzahl aller Elemente $a,a',a''\ldots a^{(\lambda+1)}$ d. i. die Zahl $\lambda+2$ ungerade oder gerade ist.

Zweitens ist erlaubt, die Ordnung der Elemente umzukehren; es wird nemlich
\[
(a,a',a''\ldots a^{(\lambda)})=(a^{(\lambda)},\ldots a'',a'\ a)
\]
Aus der Anwendung des ersten dieser Sätze auf die Grössen $g,h,k,l$ folgt
\[
gl-hk=1
\]
Die Gleichungen (4) können daher auch so dargestellt werden:
\SourceLeaf{258}
\[
\begin{aligned}
b^0&=lb^*-h\beta^*\\
\beta^0&=-kb^*+g\beta^*\\
c^0&=lc^*-h\gamma^*\\
\gamma^0&=-kc^*+g\gamma^*
\end{aligned}
\]

\Article{4}
Es sei $P$ ein gegebener Punkt auf der (nöthigenfalls verlängerten) geraden Linie, welche der erste Weg des Lichtstrahls darstellt, und $\xi,\eta,\zeta$ seine Coordinaten. Es ist also
\[
n^0\eta=\beta^0(\xi-N^0)+n^0b^0
\]
oder wenn man für $\beta^0,b^0$ die am Schluss des vorhergehenden Artikels gegebenen Ausdrücke substituirt
\[
n^0\eta=(g\beta^*-kb^*)(\xi-N^0)-n^0(h\beta^*-lb^*)
\]
folglich
\[
b^*=\frac{n^0\eta+(n^0h-g(\xi-N^0))\beta^*}{n^0l-k(\xi-N^0)}
\]
Substituirt man diesen Werth in der ersten Gleichung für den Weg des Lichtstrahls nach der letzten Brechung, nemlich in
\[
y=\frac{\beta^*}{n^*}(x-N^*)+b^*
\]
und schreibt um abzukürzen
\[
\begin{aligned}
N^*-\frac{n^0h-g(\xi-N^0)}{n^0l-k(\xi-N^0)}.n^*&=\xi^*\\
\frac{n^0\eta}{n^0l-k(\xi-N^0)}&=\eta^*
\end{aligned}
\]
so wird diese Gleichung
\[
y=\eta^*+\frac{\beta^*}{n^*}(x-\xi^*)
\]
und ganz auf dieselbe Art wird, wenn man noch
\[
\frac{n^0\zeta}{n^0l-k(\xi-N^0)}=\zeta^*
\]
schreibt, die zweite Gleichung für den Weg des Lichtstrahls nach der letzten Brechung
\SourceLeaf{259}
\[
z=\zeta^*+\frac{\gamma^*}{n^*}(x-\xi^*)
\]
Der Punkt $P^*$, dessen Coordinaten $\xi^*,\eta^*,\zeta^*$ sind, liegt also auf der (nöthigenfalls rückwärts verlängerten) geraden Linie, welche dieser letzte Weg darstellt, und zugleich ist klar, da seine Coordinaten von $\beta^0,b^0,\gamma^0,c^0$ unabhängig sind, dass er für \textit{alle} einfallenden Strahlen, die durch $P$ gehen, derselbe ist. Man kann den Punkt $P$ wie ein Object und $P^*$ als sein Bild betrachten; jenes kann aber nur dann ein reelles sein, wenn $P$ im ersten Mittel liegt, oder $\xi-N^0$ negativ ist, und eben so ist das Bild nur dann ein reelles, wenn $P^*$ in dem letzten Mittel liegt, oder $\xi^*-N^*$ positiv ist; in den entgegengesetzten Fällen sind Object oder Bild nur virtuell.

Die Punkte $P,P^*$ liegen mit der Axe der $x$ in Einer Ebene, in Entfernungen von derselben, die sich wie die Einheit und die Zahl $\frac{n^0}{n^0l-k(\xi-N^0)}$ verhalten, wobei das positive oder negative Zeichen dieser Zahl die Lage jener Punkte auf Einer Seite der Axe oder auf entgegengesetzten anzeigt. Ein System von Punkten in derselben gegen die Axe der $x$ senkrechten Ebene kann wie ein zusammengesetztes Object betrachtet werden, dessen zusammengesetztes Bild gleichfalls in Eine gegen die Axe der $x$ senkrechte Ebene fällt und dem Object ähnlich ist, so dass das Linearverhältniss der Theile durch die Zahl
\[
\frac{n^0}{n^0l-k(\xi-N^0)}=g+\frac{k}{n^*}(\xi^*-N^*)
\]
ausgedrückt wird, deren Zeichen die aufrechte oder verkehrte Lage unterscheidet.

\Article{5}
Das bisher entwickelte enthält die ganze Theorie der Veränderungen, welche der Weg der Lichtstrahlen durch Brechungen erleidet, und lässt sich leicht auch auf den Fall ausdehnen, wo mit Brechungen eine oder mehrere Reflexionen verbunden sind, was jedoch speciell hier nicht ausgeführt werden soll. Es ist aber nicht überflüssig, die Resultate in eine andere Form zu bringen, indem man sie, anstatt auf die erste und letzte Fläche oder auf die Punkte $N^0,N^*$, auf zwei andere Punkte $Q,Q^*$ bezieht. Es seien
\[
\begin{aligned}
y&=\frac{\beta^0}{n^0}(x-Q)+B\\
z&=\frac{\gamma^0}{n^0}(x-Q)+C
\end{aligned}
\]
\SourceLeaf{260}die Gleichungen für den ersten, und
\[
\begin{aligned}
y&=\frac{\beta^*}{n^*}(x-Q^*)+B^*\\
z&=\frac{\gamma^*}{n^*}(x-Q^*)+C^*
\end{aligned}
\]
die Gleichungen für den letzten Weg des Lichtstrahls, und man setze
\[
\frac{N^0-Q}{n^0}=\theta,\qquad \frac{Q^*-N^*}{n^*}=\theta^*
\]
Wir haben also
\[
\begin{aligned}
b^0&=B+\theta\beta^0,&c^0&=C+\theta\gamma^0\\
B^*&=b^*+\theta^*\beta^*,&C^*&=c^*+\theta^*\gamma^*
\end{aligned}
\]
Hieraus, verbunden mit den Gleichungen (4), folgt leicht, dass, wenn man
\[
\begin{aligned}
G&=g+\theta^*k\\
H&=h+\theta g+\theta\theta^*k+\theta^*l\\
K&=k\\
L&=l+\theta k
\end{aligned}
\]
setzt,
\[
\begin{aligned}
B^*&=GB+H\beta^0,&C^*&=GC+H\gamma^0\\
\beta^*&=KB+L\beta^0,&\gamma^*&=KC+L\gamma^0
\end{aligned}
\]
sein wird. Die Coëfficienten $G,H,K,L$, welche auf diese Weise an die Stelle von $g,h,k,l$ treten, geben auch die Gleichung
\[
GL-HK=1
\]

\Article{6}
Der Zweck der Einführung anderer Punkte, um die Lage des einfallenden und des ausfahrenden Strahls darauf zu beziehen, geht dahin, eine einfachere Abhängigkeit der letztern von der erstern darzubieten, und dazu sind vorzugsweise zwei Paare von Punkten geeignet, die mit $E,E^*$ und $F,F^*$ bezeichnet werden sollen. Die Werthe der dabei in Betracht kommenden Grössen werden sich bequem in einer tabellarischen Form übersehen lassen.
\SourceLeaf{261}
\[
\renewcommand{\arraystretch}{1.8}
\begin{array}{l|l|l}
&\multicolumn{1}{c|}{\mathrm{I}}&\multicolumn{1}{c}{\mathrm{II}}\\\cline{2-3}
\theta&\dfrac{1-l}{k}&-\dfrac{l}{k}\\
\theta^*&\dfrac{1-g}{k}&-\dfrac{g}{k}\\
Q&E=N^0-\dfrac{n^0(1-l)}{k}&F=N^0+\dfrac{n^0l}{k}=E+\dfrac{n^0}{k}\\
Q^*&E^*=N^*+\dfrac{n^*(1-g)}{k}&F^*=N^*-\dfrac{n^*g}{k}=E^*-\dfrac{n^*}{k}\\
G&1&0\\
H&0&-\dfrac1k\\
K&k&k\\
L&1&0
\end{array}
\]
Das Resultat ist also, dass, wenn die Gleichungen für den einfallenden Strahl in die Form gebracht werden
\[
\begin{aligned}
y&=\frac{\beta^0}{n^0}(x-E)+B\\
z&=\frac{\gamma^0}{n^0}(x-E)+C
\end{aligned}
\]
oder in folgende (wo wir die constanten Theile zur Unterscheidung von der ersten Form mit Accenten bezeichnen)
\[
\begin{aligned}
y&=\frac{\beta^0}{n^0}(x-F)+B'\\
z&=\frac{\gamma^0}{n^0}(x-F)+C'
\end{aligned}
\]
die Gleichungen für den ausfahrenden Strahl sein werden
\[
\begin{aligned}
y&=\frac{\beta^0+kB}{n^*}.(x-E^*)+B\\
z&=\frac{\gamma^0+kC}{n^*}.(x-E^*)+C
\end{aligned}
\]
oder
\[
\begin{aligned}
y&=\frac{kB'}{n^*}.(x-F^*)-\frac{\beta^0}k\\
z&=\frac{kC'}{n^*}.(x-F^*)-\frac{\gamma^0}k
\end{aligned}
\]
